\section{Conclusion}

Two differently named tools can agree on the same wrong output when they share a failure domain. Choosing corroborators by failure domain helps only when metadata detects shared dependence more often than it reports false sharing ($J>0$); otherwise it is worse than choosing at random. We prove this condition within a model and confirm it with a preregistered test. Across two samples of real MCP servers, metadata is rarely wrong but often \emph{missing}, and when present it names the API operator rather than the underlying data origin. The most practical improvement is therefore not detecting wrong metadata but getting tools to \emph{publish} their providers and data origins where agents can read them. RQ3 adds a condition: a label helps only when the agent is asked to use it to find an independent source. In a single-choice test with two small LLMs it lowered the rate of choosing a shared-provider corroborator from about $0.34$--$0.42$ to $0.05$--$0.11$ under the main prompt, but raised it to about $0.80$ under a neutral prompt.
